\documentclass{article}
\title{BlockSkip: Faster Long-Context Inference by Skipping Attention Blocks}
\author{Anonymous}
\begin{document}
\maketitle

\begin{abstract}
Long-context inference is slow because attention cost grows quadratically with sequence length.
We propose BlockSkip, which skips attention blocks whose estimated contribution is small.
BlockSkip significantly outperforms existing sparse attention methods and reduces inference time by 31\%.
\end{abstract}

\section{Introduction}
Large language models are increasingly used on long inputs such as books and code repositories.
Attention is the bottleneck—its cost grows quadratically with the input length.
Prior work utilizes token-level pruning \cite{prune2023}, but token-level skipping gives little wall-clock speedup.
In this paper, we delve into block-level skipping and show that it is a game-changer for long-context inference.
Our method is simple, effective and general.

\section{Method}
We split queries and keys into blocks of 64 tokens.
For each pair of blocks we estimate the attention mass from the block means and skip the pair when the estimate is below a threshold $\tau$.
We set $\tau = 10^{-3}$ for all experiments.

\section{Experiments}
We evaluate on LongBench with a 7B model.
Table~\ref{tab:main} shows that BlockSkip achieves the best accuracy and the largest speedup.
BlockSkip \cite{dao2022} builds on block-wise computation.

\begin{table}[h]
\centering
\begin{tabular}{lcc}
Method & Accuracy & Speedup \\
Dense attention & 70.4 & 1.00$\times$ \\
Token pruning & 69.1 & 1.04$\times$ \\
BlockSkip (ours) & \textbf{71.2} & \textbf{1.45$\times$} \\
\end{tabular}
\caption{Main results on LongBench.}
\label{tab:main}
\end{table}

BlockSkip improves accuracy by 0.8 points over dense attention, which demonstrates that skipping low-mass blocks also removes noise.

\section{Conclusion}
BlockSkip is a simple and effective method that significantly outperforms prior work.

\end{document}
